\chapter{Dynamical mean-field theory of the relative model}
\label{app:dmft}
\providecommand{\Dz}{\mathcal{D}z}
\providecommand{\E}{\mathbb{E}}

This note derives a dynamical mean-field theory (DMFT) for the scale-invariant model of
the growing-economy chapter, in which the interactions act on the average firm rather
than on absolute abundances. The model turns out to be a disordered replicator equation,
so its DMFT is the heterogeneous single-site process already used to locate $\mu_c$
\citep{aguirrelopez}, with two small additions that the scale invariance forces. The
reduction to a replicator and the resulting effective single-site process and its
self-consistency conditions are set out below.

\section{The model and its reduction to a replicator}
\label{sec:dmft-model}

We take the relative GLV in the competition convention used in the simulations,
\begin{equation}
    \dot x_i = x_i\left[\,1 - \frac{x_i}{m} - \frac{(\alpha x)_i}{m}\,\right],
    \qquad m = \frac1N\sum_j x_j = \frac MN ,
    \label{eq:dmft-glv}
\end{equation}
with couplings on a graph of mean degree $C$,
\begin{equation}
    \alpha_{ij} = A_{ij}\Big(-\frac{\mu}{C} + \frac{\sigma}{\sqrt C}\,z_{ij}\Big),\qquad
    z_{ij}\sim\mathcal N(0,1),\qquad \overline{z_{ij}z_{ji}} = \gamma,\qquad \alpha_{ii}=0 .
\end{equation}
Both the self-limitation $x_i/m$ and the interaction $(\alpha x)_i/m$ are divided by the
population mean $m$. The right-hand side of \eqref{eq:dmft-glv} is then homogeneous of
degree one in $x$, so the bracket $F_i := 1 - x_i/m - (\alpha x)_i/m$ depends on $x$ only
through ratios. The overall scale is a free direction that grows or decays exponentially
without the finite-time blow-up of the absolute model.

Introduce the relative abundances $y_i = x_i/m$, which satisfy $\langle y\rangle :=
\frac1N\sum_i y_i = 1$ by construction, and let $M = \sum_j x_j$ be the total. The
aggregate grows at the population-averaged fitness,
\begin{equation}
    g(t) := \frac{d\ln M}{dt}
          = \frac1M\sum_j \dot x_j = \sum_j \frac{x_j}{M}\,F_j
          = \frac1N\sum_j y_j F_j = \langle yF\rangle .
    \label{eq:dmft-g}
\end{equation}
Subtracting this common scale leaves a closed equation for the relative abundances,
\begin{equation}
    \dot y_i = y_i\big(F_i - g(t)\big),\qquad
    F_i = 1 - y_i - (\alpha y)_i,\qquad
    g(t) = \langle yF\rangle ,
    \label{eq:dmft-replicator}
\end{equation}
which is a replicator equation, the simplex projection of \eqref{eq:dmft-glv} in the sense
of Hofbauer's correspondence \citep{hofbauer81}. The aggregate is recovered from
$M(t) = M(0)\exp\!\int_0^t g$, and a positive stationary $g$ is a steadily growing
economy.

\paragraph{What the $-g(t)$ term is.}
The drift $-g(t)$ in \eqref{eq:dmft-replicator} is the mean-fitness subtraction of the
replicator, and it is the term that turns the equation for absolute size into one for
relative size. From \eqref{eq:dmft-g}, $g(t) = \langle yF\rangle = \sum_j w_j F_j$ is the
share-weighted average fitness of the whole economy, with shares $w_j = y_j/N = x_j/M$,
and it equals the growth rate of the aggregate, $g(t) = d\ln M/dt$. Every firm feels the
same $g(t)$: it is a global field, not something private to firm $i$. A firm therefore
gains relative size, $\dot y_i > 0$, only when its own fitness $F_i$ exceeds the
economy-wide average $g(t)$, so $-g(t)$ encodes the requirement to beat the field in
order to gain ground. Its role is to enforce the constraint $\langle y\rangle = 1$: one
checks from \eqref{eq:dmft-replicator} that $\frac{d}{dt}\langle y\rangle = \langle
yF\rangle - g\langle y\rangle = g\,(1 - \langle y\rangle)$, so $\langle y\rangle = 1$ is
preserved, and $g(t) = \langle yF\rangle$ is precisely the value that keeps it preserved.
Because it is an average over the whole population, $g(t)$ self-averages: its fluctuations
are of order $1/\sqrt N$ and vanish in the thermodynamic limit, so it enters the mean-field
description as a deterministic common drift rather than as a source of noise.

\section{The local field and the effective single-site process}
\label{sec:dmft-cavity}

We take the high-connectivity limit and treat the regular graph, so that
$k_i \equiv C$. Split the interaction felt by node $i$ into its mean and
its fluctuation,
\begin{equation}
    (\alpha y)_i
    = -\frac{\mu}{C}\sum_{j\in\partial i} y_j + \frac{\sigma}{\sqrt C}\sum_{j\in\partial i} z_{ij} y_j
    \;\longrightarrow\; -\mu M_1(t) + h_i(t),
\end{equation}
where the neighbourhood average of the mean part tends to the global one, $-\mu M_1$ with
$M_1 = \langle y\rangle$, and $h_i$ collects the disorder. The dynamical cavity argument
\citep{bunin2017,galla2018} proceeds by removing node $i$: the cavity trajectories
$y_j^{(i)}$ are independent of the couplings $\{z_{ij}, z_{ji}\}$, and reinstating $i$
perturbs each neighbour by linear response. The field that $i$ exerts on $j$ has
fluctuating part $-\frac{\sigma}{\sqrt C} z_{ji} y_i$, so with the single-site response
$G_j(t,s) = \delta y_j(t)/\delta\zeta_j(s)$,
\begin{equation}
    h_i(t)
    = \frac{\sigma}{\sqrt C}\sum_{j\in\partial i} z_{ij}\,y_j^{(i)}(t)
    \;-\; \frac{\sigma^2}{C}\sum_{j\in\partial i} z_{ij}z_{ji}\!\int_0^t\! ds\,G_j(t,s)\,y_i(s).
\end{equation}
The first sum is a sum of $C$ independent zero-mean terms; by the central limit theorem it
becomes a Gaussian noise $\xi_i(t)$ whose covariance is the population autocorrelation,
\begin{equation}
    \langle\xi_i(t)\xi_i(s)\rangle
    = \frac{\sigma^2}{C}\sum_{j\in\partial i} y_j(t)y_j(s)
    = \sigma^2\,C(t,s),\qquad C(t,s):=\langle y(t)y(s)\rangle .
\end{equation}
The second sum uses $\overline{z_{ij}z_{ji}} = \gamma$ and
$\frac1C\sum_{j\in\partial i} G_j \to G := \langle G_j\rangle$, giving the retarded Onsager
reaction $-\gamma\sigma^2\!\int_0^t G(t,s)\,y_i(s)\,ds$. Substituting into $F_i = 1 - y_i +
\mu M_1 - h_i$ and using \eqref{eq:dmft-replicator} produces the effective single-site
process
\begin{equation}
    \boxed{\;
    \dot y = y\left[\,1 - y + \mu M_1(t) - g(t)
        + \sigma\,\eta(t) + \gamma\sigma^2\!\int_0^t\! G(t,s)\,y(s)\,ds\,\right],
    \;}
    \label{eq:dmft-singlesite}
\end{equation}
with $\eta$ a unit-strength Gaussian noise of covariance $\langle\eta(t)\eta(s)\rangle =
C(t,s)$. Equation~\eqref{eq:dmft-singlesite} is the heterogeneous effective process of
\citet{aguirrelopez} with carrying capacity $K=1$, the competition sign on the mean field,
and the two structural additions forced by the scale invariance: the common drift $-g(t)$
of Section~\ref{sec:dmft-model}, and the constraint that the mean abundance is pinned,
$M_1\equiv 1$, rather than left free.

\section{Self-consistency}
\label{sec:dmft-selfconsistency}

The disorder average is now carried by four functionals of the single-site law of
\eqref{eq:dmft-singlesite}:
\begin{equation}
    M_1(t) = \E[y(t)] = 1,\qquad
    C(t,s) = \E[y(t)y(s)],\qquad
    G(t,s) = \E\!\left[\frac{\delta y(t)}{\delta\zeta(s)}\right]_{\zeta=0},
\end{equation}
\begin{equation}
    g(t) = \E[y(t)F(t)],\qquad
    F = 1 - y + \mu M_1 + \sigma\eta + \gamma\sigma^2\!\int_0^t G\,y .
    \label{eq:dmft-closures}
\end{equation}
The first is the normalisation, the middle two are the standard correlation and response
that close the noise and the memory, and the last fixes the common drift. The closure for
$g$ is forced rather than chosen: $\dot M_1 = \E[yF] - g M_1$, and demanding $\dot M_1 = 0$
at $M_1 = 1$ gives $g(t) = \E[yF]$. This has the structure of a solvable random-replicator
model, here carrying the GLV self-limitation as well. Solving these self-consistency
conditions is the standard replicator mean-field problem \citep{galla2018,roy2020}; in the
main text we take the freezing threshold and the firm-growth statistics directly from
simulation instead.

